\section{Methodology}
\label{sec:method}

% High-level overview
Our approach, [method name], consists of three main components:
(1) [Component A], (2) [Component B], and (3) [Component C].
Figure \ref{fig:overview} illustrates the overall architecture.

\begin{figure*}[t]
\centering
\includegraphics[width=0.8\linewidth]{figures/overview.pdf}
\caption{Overview of our methodology.}
\label{fig:overview}
\end{figure*}

% Detailed description
\subsection{Component A}
\label{sec:component_a}

Component A addresses [specific problem]. The key insight is [insight].

\paragraph{Mathematical Formulation}
Formally, we define [definition]. Given [inputs], we compute [output] using:
\begin{equation}
\label{eq:formula}
f(x) = \sum_{i=1}^{n} w_i \cdot \phi_i(x) + b
\end{equation}
where $w_i$ are [description], $\phi_i$ are [description], and $b$ is [description].

\paragraph{Algorithm}
Algorithm \ref{alg:component_a} describes the computation process.

\begin{algorithm}[t]
\caption{Component A Algorithm}
\label{alg:component_a}
\begin{algorithmic}[1]
\STATE {\bf Input:} $X, Y, \epsilon$ 
\STATE {\bf Output:} $\theta$ 
\STATE Initialize $\theta$ randomly
\FOR{iteration = 1 to max_iterations}
    \STATE Compute gradient: $g \leftarrow \nabla_\theta J(\theta)$
    \STATE Update: $\theta \leftarrow \theta - \alpha \cdot g$
    \STATE Check convergence: if $|g| < \epsilon$, break
\ENDFOR
\RETURN $\theta$
\end{algorithmic}
\end{algorithm}

\subsection{Component B}
\label{sec:component_b}

Component B [description]. The implementation uses [technology/algorithm].

\paragraph{Complexity Analysis}
The time complexity of Component B is O([complexity]), and space complexity is O([complexity]).

\subsection{Component C}
\label{sec:component_c}

\paragraph{Implementation Details}
We implemented Component C using [language/framework].
Key optimizations include:
\begin{itemize}
    \item [Optimization 1]: [description]
    \item [Optimization 2]: [description]
    \item [Optimization 3]: [description]
\end{itemize}

% Theoretical properties
\subsection{Theoretical Properties}

\begin{theorem}
\label{thm:convergence}
Under [assumptions], Algorithm \ref{alg:component_a} converges to a local minimum at rate O($1/t$).
\end{theorem}

\begin{proof}
[Proof here]
\end{proof}
